% 问题三子问题二、三独立审阅入口。
% 组织方式参考用户提供的数学建模竞赛论文：封面信息、摘要、关键词、目录、分问题分析与建模求解。
% 当前文件仅用于独立编译审阅，不接入 paper/main.tex。

\documentclass[UTF8,a4paper,12pt]{ctexart}
\usepackage{amsmath,amssymb,booktabs,geometry,hyperref,longtable}
\geometry{left=2.7cm,right=2.7cm,top=2.6cm,bottom=2.6cm}
\hypersetup{colorlinks=true,linkcolor=black,citecolor=black,urlcolor=black}
\setlength{\parindent}{2em}
\setlength{\parskip}{0.25em}
\sloppy
\emergencystretch=3em
\renewcommand{\contentsname}{目录}
\renewcommand{\figurename}{图}
\renewcommand{\tablename}{表}
\begin{document}

\begin{titlepage}
  \centering
  \vspace*{2.0cm}
  {\heiti\zihao{2}问题三条件资源配置分析\par}
  \vspace{0.8cm}
  {\heiti\zihao{3}子问题二与子问题三独立审阅稿\par}
  \vfill
  {\zihao{4}预算约束下的原生质量优化与配比接口敏感性\par}
  \vspace{1.0cm}
  {\zihao{4}版本：Q3-CONDITIONAL-INTERFACE-v1.0.0\par}
  \vspace{2.0cm}
  {\zihao{4}DRAFT / REVIEW REQUIRED\par}
\end{titlepage}

\pagenumbering{roman}
\section*{摘要}
本文按照问题三的条件资源配置口径，分别讨论 B6/B7 原生质量变量
$Q_B$ 下的 $N$--$D$--$Q_B$ 优化，以及 A 侧观测配比通过假设
$Q_A\to Q_B$ 接口进入 Q3 的条件面板。由于 A、B 数据缺少共同连接键，
且 $Q_A$ 与 $Q_B$ 尚未统一标度，配比结果只作为条件情景和接口敏感性结果，
不解释为跨来源联合最优。所有数值均受 B1/B6 观测范围、给定预算、上下文
和成本函数假设约束。

\noindent\textbf{关键词：} 条件资源配置；广义标度律；原生质量变量；
多起点 SLSQP；配比接口；敏感性分析

\clearpage
\tableofcontents
\clearpage
\pagenumbering{arabic}

\setcounter{section}{5}
\section{问题三的求解：预算约束下的条件资源配置}
\label{sec:q3-standalone}
问题三的总体输入包括 B1 的 $N,D$ 标量、B6/B7 的原生 $Q_B$ 以及
A 侧有限个观测配比候选。本文将这些输入拆分为可识别的 B 侧条件优化
和明确标注的 A--B 假设接口，不在缺少共同键的情况下拼接四量联合样本。

\setcounter{subsection}{4}
% 问题三子问题二候选正文。
% 仅使用 B6/B7 原生 Q_B 条件模型及已登记的 Q3 优化运行；
% 当前稿件待科学评审，不直接修改 paper/main.tex。

\subsection{问题三子问题二：M1 的规模--数据量--质量条件优化}
\label{sec:q3-subproblem2}

\subsubsection{问题分析与变量界定}

子问题二在问题二 M1 的基础上，将质量资源作为第三个决策变量。本文记该变量为
$Q_B$，它是 B6/B7 数据表中的原生 \texttt{Q\_score}，不是问题一得到的
A 侧质量变量 $Q_A$。因此，$Q_B$ 只在 B6/B7 条件模型内部解释为质量分数，
不能与 $Q_A$ 直接相等，也不能把两个量尺合并成一个已验证的质量变量。

M1 的三个决策变量为模型规模 $N$、训练数据量 $D$ 和原生质量分数 $Q_B$。
其中 $N$ 和 $D$ 分别以十亿参数和十亿 tokens 计，$Q_B$ 的观测范围为
$[0.1,1.0]$。$N$ 和 $D$ 的有界范围取自 B1 样本：
\begin{equation}
  N\in[0.070542,11.965825],\qquad
  D\in[0.134,299.893].
  \label{eq:q3-sp2-support}
\end{equation}
这些区间表示当前数据的观测适用范围，不表示对更大模型或更多训练数据的物理定律外推许可。

\subsubsection{M1 目标函数与质量成本}

冻结的 B6 M1 条件损失为
\begin{equation}
  \widehat L_{\mathrm{M1}}(N,D,Q_B)
  =E+A N^{-\alpha}+B D^{-\beta}+G(1-Q_B).
  \label{eq:q3-sp2-m1-loss}
\end{equation}
本轮使用的参数为
\begin{equation}
\begin{aligned}
E&=1.4892660413, & A&=0.5401729883, & B&=1.3167948616,\\
G&=0.3622474350, & \alpha&=0.2790539239, &
\beta&=0.2828029135.
\end{aligned}
\label{eq:q3-sp2-m1-parameters}
\end{equation}

给定上下文长度 $L_{\mathrm{ctx}}$ 时，基础规模成本写为
\begin{equation}
  C_{\mathrm{base}}
  =10^{18}ND\left(6+\eta L_{\mathrm{ctx}}\right),
  \qquad \eta=2\times10^{-4}.
  \label{eq:q3-sp2-base-cost}
\end{equation}
相对于基准质量 $Q_0=0.1$ 的质量增量成本写为
\begin{equation}
  C_Q=10^9D\left[g(Q_B)-g(Q_0)\right].
  \label{eq:q3-sp2-quality-cost}
\end{equation}
本文比较三类成本形状：
\begin{equation}
\begin{aligned}
g_{\mathrm{exp}}(Q)&=10^7\exp(6Q),\\
g_{\mathrm{pow}}(Q)&=5\times10^9Q^4,\\
g_{\mathrm{log}}(Q)&=2\times10^9\log(1+10Q).
\end{aligned}
\label{eq:q3-sp2-cost-families}
\end{equation}
三类 $g$ 用于敏感性比较，不被解释为已经由真实训练账单验证的成本定律。
在预算 $C$ 下，优化问题为
\begin{equation}
\begin{aligned}
\min_{N,D,Q_B}\quad&
\widehat L_{\mathrm{M1}}(N,D,Q_B),\\
\mathrm{s.t.}\quad&
C_{\mathrm{base}}(N,D;L_{\mathrm{ctx}})
+C_Q(D,Q_B)\le C,\\
&N\in[0.070542,11.965825],\\
&D\in[0.134,299.893],\\
&Q_B\in[0.1,1.0].
\end{aligned}
\label{eq:q3-sp2-opt}
\end{equation}

质量成本占比定义为
\begin{equation}
  \rho_Q=\frac{C_Q}{C_{\mathrm{base}}+C_Q}.
  \label{eq:q3-sp2-quality-share}
\end{equation}
$\rho_Q$ 是当前成本假设下的模型内比例，不是现实工程账单中的实测占比。

\subsubsection{求解方法}

为保持正值约束，$N$ 和 $D$ 采用对数参数化，优化变量为
$(\log N,\log D,Q_B)$。预算约束按预算值缩放后送入 SLSQP，以避免
$10^{19}$--$10^{24}$ 量级的原始约束造成有限差分尺度失衡。

每个情景使用六个确定性初始点，覆盖低值、几何中点、高值以及不同的
$Q_B$ 初始值。多起点的目的在于检查局部求解是否受初始点影响，并比较
可行性、收敛状态、目标函数值和边界状态。最终优先保留优化器报告收敛且预算可行的解。
当前不采用遗传算法或粒子群算法，原因是决策维度低、目标函数连续光滑、约束边界明确，
解析结构与多起点局部优化更便于逐情景审计。

求解后将结果分为三类：$Q_B=0.1$ 或 $Q_B=1.0$ 的端点状态；
$0.1<Q_B<1.0$ 且 $N,D$ 未触及边界的内点状态；以及 $N$、$D$ 或 $Q_B$
触及观测范围边界的边界状态。

\subsubsection{\texorpdfstring{原生 $Q_B$ 情景结果}{原生 Q\_B 情景结果}}

本轮共生成 $3\times5\times3=45$ 个原生 $Q_B$ 情景，覆盖三类质量成本、
五档上下文
\[
L_{\mathrm{ctx}}\in\{2048,4096,8192,32768,131072\}
\]
和三档预算
\[
C\in\{10^{19},10^{22},10^{24}\}\ \mathrm{FLOPs}.
\]
每行结果报告 $Q_B^*$、$N^*$、$D^*$、$L^*$、质量成本占比、预算活跃状态、
边界状态和多起点收敛信息。

按上下文取平均后，低预算 $C=10^{19}$ 下，指数型、幂函数型和对数渐进型
的最优质量分别约为 $0.7243$、$0.6103$ 和 $1.0000$。预算提高到
$10^{22}$ 或 $10^{24}$ 后，当前参数下三类成本通常均将 $Q_B^*$ 推到
$1.0000$。这只说明在当前 M1 和成本假设下质量增量成本相对 $N,D$ 成本较小，
不表示真实训练质量已经饱和，也不表示 $Q_B=1$ 是经验验证的唯一最优质量。

在指数成本、$L_{\mathrm{ctx}}=8192$、$C=10^{22}$ 的代表性情景中，质量成本
约占总成本的 $5.3\%$--$7.4\%$。高预算情景中若同时触及
$N=11.965825$、$D=299.893$ 和 $Q_B=1.0$，则应将其标记为观测范围边界解；
名义预算的剩余部分不能用来证明支持范围外仍有有效收益。

\begin{table}[htbp]
  \centering
  \caption{M1 原生质量条件情景的摘要}
  \label{tab:q3-sp2-summary}
  \small
  \begin{tabular}{lccc}
    \toprule
    质量成本族 & $C=10^{19}$ 的平均 $Q_B^*$ &
    $C=10^{22}$ 的典型状态 & 主要边界 \\
    \midrule
    指数型 & 0.7243 & 接近 $1.0$ & 低预算内点、高预算边界 \\
    幂函数型 & 0.6103 & 接近 $1.0$ & 低预算内点、高预算边界 \\
    对数渐进型 & 1.0000 & $1.0$ & 质量上界较早活跃 \\
    \bottomrule
  \end{tabular}
\end{table}

\subsubsection{上下文成本结构}

由式~\eqref{eq:q3-sp2-base-cost}，基础成本中的注意力项与训练项之比为
\begin{equation}
  \frac{C_{\mathrm{attn}}}{C_{\mathrm{train}}}
  =\frac{\eta L_{\mathrm{ctx}}}{6}
  =\frac{L_{\mathrm{ctx}}}{30000}.
  \label{eq:q3-sp2-context-ratio}
\end{equation}
因此，$L_{\mathrm{ctx}}=30000$ 是当前成本结构中的代数临界点：
低于该值时注意力项小于训练项，高于该值时注意力项大于训练项。
它不是已经由任务效果验证的上下文性能临界点，也不是 $Q_B$ 的经验转折点。

\subsubsection{数值复核与稳健性}

本轮复核包括以下内容：
\begin{enumerate}
  \item 45 行原生 $Q_B$ 输出及参数折叠传播结果均为有限值；
  \item 原生质量情景最大相对预算误差约为 $2.1\times10^{-12}$；
  \item 质量成本稳健性运行的最大相对预算误差约为
  $9.7\times10^{-12}$，参数折叠运行的 M1 最大相对误差约为
  $9.01\times10^{-12}$；
  \item $N,D$ 未超过 B1 观测范围，$Q_B$ 未超过 B6 原生范围；
  \item 六个起点的成功状态、目标值差异和边界标志均被保留；
  \item 原生情景表 SHA256 为
  {\scriptsize\texttt{77171f327c40d0bbd2559665b5aa7b6fdc945b5052301dbd1c73296289330378}}；
  \item $Q_0\in\{0.1,0.3,0.5\}$ 的扫描显示，低预算下指数型和幂函数型
  对 $Q_0$ 有可见敏感性，对数渐进型在当前参数下系统性地推向质量上界。
\end{enumerate}

\subsubsection{结果讨论与小结}

在 B1 的 $N,D$ 观测范围、B6/B7 原生 $Q_B$ 范围以及给定预算、上下文和
质量成本假设下，可以报告质量--规模--数据量的条件替代关系、成本族造成的
资源配置差异、上下文改变后的训练--注意力成本转移，以及参数折叠和边界检查
支持的稳健范围。

不能据此声称 A--B 联合 Loss 已经验证，不能声称 $Q_A$ 与 $Q_B$ 已经建立稳定
实证关系，不能声称 $Q_B=1$ 是真实训练中的唯一最优质量，也不能把支持范围外的
结果解释为现实中的全局最优。以上结论均属于冻结 M1 参数和成本假设下的条件情景。

本子问题的证据运行包括：
\begin{itemize}
  \item \texttt{q3-q-conditional-20260925-r01}；
  \item \texttt{q3-q-robustness-20260925-r01}；
  \item \texttt{q3-optimization-robustness-20260926-r01}；
  \item \texttt{q3-answer-package-20260926-r01}。
\end{itemize}


\setcounter{subsection}{5}
% 问题三子问题三候选正文。
% A 侧配比只作为条件标签；Q_A 到 Q_B 的传递是未经验证的假设接口。

\subsection{问题三子问题三：M2/P 配比条件面板与 Q3 接口敏感性}
\label{sec:q3-subproblem3}

\subsubsection{配比进入方式}

子问题三使用 A4/A5 中已经观测到的六个配比候选。记配比向量为
$p=(p_1,\ldots,p_J)$，满足
\begin{equation}
  p_j\geq0,\qquad \sum_{j=1}^{J}p_j=1.
\end{equation}
A 侧候选质量分数按
\begin{equation}
  Q_A(p)=q^\top p=\sum_{j=1}^{J}q_jp_j
  \label{eq:q3-sp3-qa}
\end{equation}
计算。$Q_A(p)$ 只作为 A 侧配比条件标签，不直接加到 B 侧 M1 的
Loss 中，也不替换 M1 中的 $Q_B$。

本节保留两个层次。第一层是六个观测配比候选与 Q3 的 15 个预算--上下文
$N$--$D$ 情景组成的 90 行条件面板；第二层是在明确标注的
$Q_A\to Q_B$ 假设接口下，将固定的假设 $Q_B$ 传给 M1 求解器，形成
接口敏感性结果。90 行和 1350 行均表示实验覆盖量，不是新增观测样本量。

\subsubsection{假设质量接口}

由于 A 侧 $Q_A$ 与 B6/B7 原生 $Q_B$ 没有共同量尺和样本级连接键，本文不把
二者直接相等。为进行可计算的接口压力测试，使用登记的区间仿射校准：
\begin{equation}
  Q_B^{\mathrm{hyp}}(p)
  =0.1+0.9\,\mathrm{clip}
  \left(
  \frac{Q_A(p)-q_{\min}}{q_{\max}-q_{\min}},0,1
  \right),
  \label{eq:q3-sp3-hyp-map}
\end{equation}
其中
\begin{equation}
  q_{\min}=51.7014075324,\qquad
  q_{\max}=64.3357099536.
  \label{eq:q3-sp3-map-range}
\end{equation}
该变换仅把 Q1 候选的 $0$--$100$ 量尺压缩到 B6 的 $[0.1,1.0]$ 计算区间，
没有使用 A--B 配对样本估计校准参数。

正文必须保留三条限制：
\begin{enumerate}
  \item 这不是 $Q_A,Q_B$ 的实证关系；
  \item 在没有共同键和标度验证的情况下，不能识别 $G_{\mathrm{bridge}}$；
  \item 映射结果是反事实条件情景，不是新增观测样本。
\end{enumerate}

\subsubsection{映射情景与计算设计}

本轮并列保留五种 Q1 侧构造：部分映射原值
\texttt{M2\_partial\_raw}、部分映射重归一化
\texttt{M2\_partial\_renorm}，以及对未映射质量使用低、中、高端点的
\texttt{M2\_interval\_low}、\texttt{M2\_interval\_mid} 和
\texttt{M2\_interval\_high}。五种构造仍在 Q1 的 $0$--$100$ 量尺中保存，
只有进入 Q3 接口压力测试时才使用式~\eqref{eq:q3-sp3-hyp-map}。

接口敏感性矩阵为
\begin{equation}
  6\times5\times3\times5\times3=1350
\end{equation}
行，分别对应六个候选、五种映射、三类质量成本、五档上下文和三档预算。
每行报告 $Q_A$、假设 $Q_B$、$N^*$、$D^*$、$L^*$、质量成本占比、
预算活跃状态、边界状态和支持范围状态。M0 与 B6 原生 M1 只作为独立参考，
不参与参数平均，也不与接口行合并拟合。

\subsubsection{配比条件面板结果}

六个候选行号为 46、97、170、186、270 和 323。90 行面板用于并列考察
观测配比候选与问题三规模资源情景之间的条件组合。它不表示 A 侧每一行与
B1 的某一行存在对应关系，也不表示 $p$ 已经进入 $N$--$D$ 目标函数。

观测标准化 Loss 最低的候选是 170，而跨模型平均预测排名第一的是 46。
排序差异说明当前离散候选面板不能推出稳定的唯一最优配方，配比结果应与
Q3 的 $N$--$D$ 资源解分列呈现。

\subsubsection{Q2--Q3 接口敏感性结果}

在指数型质量成本、$L_{\mathrm{ctx}}=2048$、$C=10^{22}$ FLOPs 的代表性情景下，
按六个候选取中位数，接口构造得到表~\ref{tab:q3-sp3-representative}。

\begin{table}[htbp]
  \centering
  \caption{代表性预算--上下文下的接口情景中位数}
  \label{tab:q3-sp3-representative}
  \small
  \begin{tabular}{lrrrrrr}
    \toprule
    映射情景 & $Q_A$ & 假设 $Q_B$ & $N^*$ & $D^*$ & $L^*$ & $\rho_Q$ \\
    \midrule
    partial raw & 33.39 & 0.100 & 8.094 & 192.756 & 2.414 & 0.000 \\
    interval low & 57.21 & 0.493 & 8.132 & 191.168 & 2.272 & 0.004 \\
    interval mid & 61.83 & 0.822 & 8.359 & 181.994 & 2.155 & 0.025 \\
    partial renorm & 62.24 & 0.851 & 8.407 & 180.129 & 2.145 & 0.029 \\
    interval high & 63.54 & 0.943 & 8.630 & 171.916 & 2.113 & 0.049 \\
    \bottomrule
  \end{tabular}
\end{table}

在这一代表性场景中，假设 $Q_B$ 增大时，$N^*$ 增大、$D^*$ 减小、
质量成本占比上升而 Loss 下降。全体 1350 行情景保持这一条件方向，
但该方向只描述冻结 M1 和给定成本函数下的资源替代，不证明真实系统中
提高 A 侧配比质量一定会提高 B6 的 $Q_B$。

六个 \texttt{partial\_raw} 候选均被压到 $Q_B=0.1$ 下界，说明“未映射质量
直接视为零贡献”的构造会把接口解推向低质量端。这是接口假设的结果，不能
解释成 B6 观测到的低质量事实。作为独立参照，原生 B6 M1 优化在同一类
预算--上下文条件下给出 $Q_B=1.0$、$N^*=8.834$、$D^*=164.915$、
$L^*=2.094$，该结果不能与 M2/P 假设映射情景合并解释。

\subsubsection{数值复核与稳健性}

接口结果复核如下：
\begin{enumerate}
  \item 六个候选的 $Q_A$、假设 $Q_B$ 及全部 1350 行 Q3 输出均为有限值；
  \item 所有假设 $Q_B$ 均位于 B6 的 $[0.1,1.0]$ 区间，$N,D$ 均位于
  B1 观测范围；
  \item 最大相对预算误差为 $9.3193240576\times10^{-12}$；
  \item \texttt{m2\_q3\_fixed\_q\_scenarios.csv} 的 SHA256 为
  {\scriptsize\texttt{a5009de0b6cc88d43b61a65dccaec79600bbbbc060f4b610b269d1be9dbb95d4}}；
  \item \texttt{m2\_q3\_interface\_summary.csv} 的 SHA256 为
  {\scriptsize\texttt{f410270c93b2fefed1af2f56463786835740354d701d802e3b503750a150407e}}；
  \item 接口检查保持 \texttt{a\_b\_row\_join=false}、
  \texttt{formal\_m3\_fit=false} 和 \texttt{b8\_used=false}。
\end{enumerate}

\subsubsection{结果讨论与主张边界}

在当前证据范围内，可以报告六个已观测配比候选的 $Q_A(p)$ 条件面板，
五种接口构造在不同成本、预算和上下文条件下对 $Q_B,N^*,D^*,L^*$ 及
质量成本占比的条件影响，以及映射构造造成的资源配置变化。

以下内容只能写成条件结论：在某个明确的 $Q_A\to Q_B$ 假设下，
某个候选配比对应的 Q3 资源解如何变化；接口选择如何改变预算前沿；
以及 M3-H 有效数据效率模型的理论反事实结果。

不能据此声称 A--B 联合 Loss 已经验证，不能声称 $Q_A$ 与 $Q_B$ 已经建立
稳定关系，不能声称 $G_{\mathrm{bridge}}$ 已被识别，不能声称六个候选中
存在现实中的唯一最优配比，也不能把 1350 行情景写成 1350 个新增观测样本。

\subsubsection{小结}

在给定的观测配比候选、预算、上下文和成本假设下，本节完成了 A 侧配比条件面板
以及 $Q_A\to Q_B$ 假设接口的敏感性分析。结果表明，不同映射会改变
$Q_B$、$N^*$、$D^*$、Loss 和质量成本占比，因此接口选择是问题三结果的重要
敏感性来源。由于 A、B 数据缺少共同连接键，且 $Q_A$ 与 $Q_B$ 尚未统一标度，
本文不把这些反事实情景解释为跨来源联合最优，也不据此选择唯一配比。

本子问题的证据运行包括：
\begin{itemize}
  \item \texttt{q3-p-conditional-20260925-r01}；
  \item \texttt{q2-q3-interface-sensitivity-20260926-r01}；
  \item \texttt{q3-answer-package-20260926-r01}。
\end{itemize}


\end{document}
